% RESEARCH_DRAFT_SHA256: f0abc89868f51a4abf5329838d8ddd0f7f1c3c24215a258a34c27725f36365f9
% Layout and source-reference wrappers expanded for the web renderer.
\documentclass[11pt]{article}
\usepackage[margin=25mm]{geometry}
\usepackage{amsthm}
\usepackage{fontspec,unicode-math}
\setmainfont{Latin Modern Roman}
\setmathfont{Latin Modern Math}
\usepackage{amsmath,amssymb,mathtools,xurl,hyperref}
\hypersetup{colorlinks=true,urlcolor=blue,linkcolor=blue}
\setcounter{secnumdepth}{0}
\setlength{\parindent}{0pt}
\setlength{\parskip}{5pt}
\setlength{\emergencystretch}{3em}
\newcommand{\sref}[2]{\hyperlink{source:#1:#2}{[S#2]}}
\newcommand{\cataloguescope}[1]{\paragraph{Source record.}#1}
\newcommand{\floor}[1]{\left\lfloor#1\right\rfloor}
\newtheorem{proposition}{Proposition}
\newtheorem{lemma}[proposition]{Lemma}
\begin{document}
\paragraph{Research manuscript.} OpenAI \texttt{gpt-6-astra}, reasoning \texttt{ultra}; 11 October 2026.
\paragraph{Preserved source body.} The complete catalogue mathematical text and sources follow with layout and source links expanded, including source qualifications and its historical preparation notes. New research begins at the section headed Research attempt.

\section{Human Alignment of Artificial Intelligence Systems}\label{statement}

\subsection{Definitions and mathematical statement}

\paragraph{Shared notation.}Standard number systems are $\mathbb N,\mathbb Z,\mathbb Q,\mathbb R,\mathbb C$. Specialized notation and index ranges follow the cited source and the definitions below.


A representation maps observations into features used for subsequent decisions. Alignment with human behavior can be evaluated only after specifying human judgments, a comparison measure, and the tasks being compared.
\paragraph{Record scope and qualifications.}
The catalogue title is retained as metadata. Named certificates, contracts, bounds or frameworks in that title are not asserted as results of the original source extract.
The source gives a broad alignment direction without a mathematical alignment criterion or quantified target.

Source-record limitation: Human-like AI behavior and uncertainty are requested without mathematically defined alignment criteria.
\paragraph{Statement recorded in the supplied source.}
How can artificial intelligence systems be aligned more closely with human behavior and uncertainty, particularly by representing the full fine-to-coarse hierarchy of human conceptual knowledge? \href{https://doi.org/10.4171/owr/2026/11}{[S2]}
\paragraph{Source status.}
The supplied catalogue labels this source record partially solved; that metadata was checked against the supplied snapshot on 2026-10-11. The label is preserved for provenance and is not an independent verification of present open status or mathematical correctness. The catalogue literature assessment is dated 2026-08-22.
\subsection{Short English statement}
Improve agreement between AI systems and human behavior and conceptual structure.
\subsection{Sources}
\begin{itemize}
\item[S1] Ulam AI and the UnsolvedMath Contributors, UnsolvedMath: A Curated Collection of Open Mathematics Problems, catalogue entry. Catalogue curation and metadata: CC BY 4.0. The source extract is retained; definitions, English summaries and qualifications are editorial additions.
\url{https://huggingface.co/datasets/ulamai/UnsolvedMath}.
\item[S2] Modern and Emerging Phenomena in Machine Learning (2026), Oberwolfach Reports, DOI 10.4171/owr/2026/11, pp. 683--736
\url{https://doi.org/10.4171/owr/2026/11}.
\item[S3] Muttenthaler et al., Aligning Machine and Human Visual Representations
\url{https://arxiv.org/abs/2409.06509}.
\item Source-specific Oberwolfach Reports licences, including CC BY-SA 4.0 where applicable, are retained; see THIRD\_PARTY\_NOTICES.md.
\end{itemize}
\paragraph{Preparation and review.}
The original source and UnsolvedMath contributors retain attribution for the source extract. Codex (GPT-6; exact serving revision not exposed) assisted with definitions, English summaries, scope notes and formatting. No new human mathematical authorship or independently verified solution is asserted. Catalogue identity checked against the complete supplied registry and current repository. Distinct upstream records, including repeated questions and broad or damaged extracts, are retained for complete coverage of the source CS category. A repeated record is not a new independent conjecture. Complete mathematical-equivalence and current-literature review has not been performed. Preparation and final review follow contribution prompts C and D, with source gaps disclosed. Structural and rendering checks do not establish a new theorem or open status.
\label{end:statement}

\section{Research attempt}
\subsection{Literature, source context and scope}
The full catalogue body, including qualifications and historical source status, is preserved above. The original workshop report was inspected in the locally cached publisher PDF text. This attempt distinguishes the source question from any explicitly restricted model introduced below. Repeated records retain their own source IDs and do not become independent conjectures. Elementary results and attributed literature results are not claimed as new original-target progress.

\subsection{Testing coarse agreement as a certificate of hierarchical alignment}
The source concerns agreement with human judgments across conceptual levels. No quantitative acceptance test is fixed. We introduce a limited probability model: a finite rooted tree of categories, with positive distributions $P,Q$ on its leaves. For each node $v$, write $P(v),Q(v)$ for the mass of its descendant leaves. At each internal node, $P(\cdot\mid v)$ and $Q(\cdot\mid v)$ are distributions over its children.
\begin{proposition}
The leaf divergence decomposes exactly as
\[
 D_{\rm KL}(P\|Q)=\sum_{v\text{ internal}} P(v)
 D_{\rm KL}(P(\cdot\mid v)\|Q(\cdot\mid v)).
\]
Agreement at a coarse level alone does not bound fine-level total variation away from one.
\end{proposition}
\begin{proof}
For a leaf $\ell$, its probability is the product of successive child-given-parent probabilities along its root-to-leaf path. Hence $\log(P(\ell)/Q(\ell))$ is the sum of the corresponding conditional log ratios. Multiply by $P(\ell)$ and sum over leaves. Interchanging the finite sums, the coefficient associated with a parent-child edge $v\to u$ is the total leaf mass $P(u)=P(v)P(u\mid v)$. Grouping edges by parent gives the identity.
For the second assertion take four leaves grouped into two pairs and let
\[
 P=(1/2-\varepsilon,\varepsilon,\varepsilon,1/2-\varepsilon),\quad
 Q=(\varepsilon,1/2-\varepsilon,1/2-\varepsilon,\varepsilon),
 \quad0<\varepsilon<1/4.
\]
Both group-mass vectors equal $(1/2,1/2)$, but their leaf total variation is $\tfrac12\sum_i|P_i-Q_i|=1-4\varepsilon$, approaching one.
\end{proof}
The attempted inference from a perfect coarse-level score to full hierarchical agreement therefore fails. The identity indicates precisely which conditional discrepancies the coarse score omits. It concerns a declared probability representation, not a claim that human conceptual knowledge literally forms this tree or that KL is the correct behavioral metric.
\subsection{Verification and remaining gaps}
At $\varepsilon=1/100$, coarse discrepancy is zero, leaf total variation is $0.96$, and direct leaf KL equals $0.96\log49$, matching the conditional decomposition. No human data or alignment intervention was evaluated. Choosing a valid hierarchy, judgment law and predictive metric remains external to the theorem.

\section{References and source attribution}
The primary Oberwolfach workshop report, \url{https://doi.org/10.4171/owr/2026/11}, was inspected on 11 October 2026. Relevant source passages: 2026/11 pp. 712--713, fine-to-coarse human conceptual alignment. Catalogue and contributor attribution remain in the original Sources section. Mathematical arguments labeled as this attempt are proved above; cited external results remain attributed to their authors and are not independently reproved or claimed as our discoveries.

\url{https://arxiv.org/abs/2409.06509}. Muttenthaler and coauthors, Aligning Machine and Human Visual Representations across Abstraction Levels. The primary abstract was opened and its task-specific representation-transfer context checked. Its empirical improvements are not claimed as findings of this attempt, which collected no human-judgment data.


\paragraph{Integration review.} Codex (GPT-6; exact serving revision not exposed) reviewed the written arguments and expanded layout and source-reference wrappers. Repository metadata and the production web renderer were checked using Python 3.13 and JavaScript with MathJax 3.2.2. This model-assisted review is not independent human proof verification; no PDF compilation or formal proof checking is asserted.

\section{Completion Estimate}
\noindent\textbf{COMPLETION: 0\%}
No substantiated new progress toward the entire original target is claimed. The worked results and failed route are delimited above; a known result, illustrative model, or auxiliary obstruction is not counted as original-target completion.

\section{Final status and attribution}
\textbf{Attempt status: unresolved.} This describes this research attempt, not a new ruling on the source catalogue's status or on external literature claims.
The mathematics and this self-audit were generated by OpenAI \texttt{gpt-6-astra}, reasoning \texttt{ultra}, on 11 October 2026. Codex (GPT-6; exact serving revision not exposed) orchestrated the batch and contributed formatting/repository review. No human mathematical authorship, independent human proof checking, or formal verification is claimed. The historical preparation paragraphs in the reproduced catalogue body concern that body only.

Contribution prompt A was applied to the complete definition and supplied source record, followed by the adversarial checks of prompt B. Local reproduction files are \texttt{batch011/checks.py}, \texttt{computations.json}, and the per-ID audit JSON. They were not public links at preparation time. Numerical checks are finite and do not prove unrestricted results. No LaTeX engine was found locally; compilation was not performed by this research agent.

\end{document}
